Four experiments were conducted: a single-rate characterization of the
front end and its error modes, a characterization of the link across data
rates, an offline evaluation of the decoder against held-out transmissions, and
a live closed-loop demonstration on hardware. All measurements were taken with
the beacon and receiver in a fixed line-of-sight geometry at approximately
0.1~m, under ordinary indoor lighting. No result reported here varies the
range, and none should be read as characterizing the link's behaviour with
distance.

\subsection{Automated Data Collection}

Labeled data were collected by driving the beacon over its network interface and
capturing the receiver's pulse stream. Because the transmitted message is set
programmatically, the ground-truth symbol sequence is known exactly: no
alignment heuristic is required, and there is consequently no risk of deriving
labels from the same duration feature the model is later trained on.

The beacon re-reads its message at the start of each transmission, so a change
takes effect only at a message boundary; the transmission in flight when the
message changes is therefore discarded. Transmissions are separated in the
captured stream by the 500~ms inter-message rest rather than by decoding them,
keeping segmentation independent of the process under evaluation. Captures whose
length departs substantially from the expected symbol count are rejected as
partial or straddling a message change.

Payloads are drawn from three families. In the forty-message set used for the
rate sweep, twenty-two messages are natural English prose sliced from a fixed
corpus on word boundaries; ten are random printable ASCII drawn over letters,
digits, punctuation and space; and eight are hand-constructed edge cases
chosen to exercise framing, including
runs of a single repeated character, alternating bit patterns, and bytes whose
set bits sit at the ends of the least-significant-bit-first sequence. The three
families are interleaved in transmission order so that payload type cannot
become confounded with time. Retaining both prose and random payloads is
deliberate: comparing them is what would separate a future language prior's
contribution from what the structural constraints already provide.

\subsection{Front-End and Error-Mode Characterization}

A single-rate dataset was collected at 28.6~bit/s to characterize the front end
and the baseline decoder's error modes with enough symbols to resolve rare
events: 437 transmissions comprising 67{,}581 symbols, of which 283 carried
natural text. Every symbol in this dataset was detected, which is what makes it
usable for measuring misclassification in isolation from symbol loss. Per-pulse
peak amplitude and peak filter response were recorded alongside duration.

\subsection{Data-Rate Sweep}

The link was characterized at five rate multipliers: 1$\times$, 2$\times$,
3$\times$, 4$\times$ and 6$\times$, corresponding to 28.6, 57.1, 85.7, 114.3
and 171.4~bit/s. Forty distinct messages were transmitted at each rate, three
captures retained per message, giving 597 labeled transmissions comprising
approximately 91{,}000 pulses. The same forty messages were used at every rate
so that differences are attributable to rate rather than payload. One capture
was discarded by the analysis for carrying a serial line truncated mid-value.

Before each block the receiver's adaptive threshold is allowed to reconverge,
and the serial buffer is flushed, so that captures made while the detector is
still adapting do not enter the dataset.

Unless stated otherwise, every figure and table below is computed from this
sweep. Where a quantity is available from more than one collection, the sweep
is the one reported, so that no two numbers in a single comparison come from
different runs.

\subsection{Offline Decoder Evaluation}

The decoder was evaluated against the threshold decoder implemented in the
receiver firmware. Two aspects of the protocol make this comparison meaningful
rather than rigged.

First, the baseline's acceptance windows are \emph{scaled with the rate}.
Leaving them at their nominal settings would cause the baseline to fail
completely above 1$\times$, which would be an artifact of the comparison rather
than a result. The question posed is whether a fitted model outperforms
hand-tuned thresholds that were themselves retuned for each operating point.

Second, train and test are split \emph{by message}, never by capture. The three
captures of a given message share a payload, so splitting at the capture level
would place the same content on both sides of the split. Thirty percent of
messages at each rate were held out, giving 36 test transmissions per rate and
35 at 171.4~bit/s. The eight parameters of Table~\ref{tab:params} were fitted
on the remaining 28 messages, independently at each rate.

The split is deterministic rather than random: messages are sorted and the
first twelve held out. This has a consequence that must be stated, because it
bears on how the offline numbers should be read. Sorting by character code
places punctuation, digits and uppercase ahead of lowercase prose, so the
held-out set contains \emph{no} natural-text messages at all. All twelve are
random or hand-constructed payloads, while 22 of the 28 training messages are
English. Two things follow. The comparison between payload families motivated
above cannot be performed on this split, and is left to future work. And the
offline results are measured on payloads with no linguistic structure and a
deliberate concentration of framing edge cases, which if anything understates
the decoder relative to the traffic a deployed link would carry. Held-out
messages also run longer than average, 174 symbols against 153 across the
sweep, which matters in Section~\ref{results}.

\subsubsection{Metrics}

Two metrics are reported, and both are defined over \emph{received
transmissions} rather than over distinct payloads: one transmission is one
complete emission of a message, so a message sent three times contributes three
independent trials. Naming the unit this way is deliberate. The term
\emph{packet} would be misleading, because the link carries nothing of the kind
in the networking sense, having no header, no length field and no checksum;
and \emph{message} alone would be ambiguous, because the same forty messages
are transmitted repeatedly and at every rate. The unit of delivery is one
newline-terminated message, and the unit of measurement is one received
transmission of it.

Let transmission $i$ carry the byte sequence
$\mathbf{b}^{(i)}=(b_1,\dots,b_{L_i})$, comprising its payload characters and
the newline the beacon appends as terminator, and let $\hat{\mathbf{b}}^{(i)}$
be what the decoder returns for it. Over $N$ evaluated transmissions, the
transmission success rate is
\begin{equation}
\mathrm{TSR} = \frac{1}{N}\sum_{i=1}^{N}
  \mathbf{1}\!\left[\hat{\mathbf{b}}^{(i)} = \mathbf{b}^{(i)}\right],
\label{eq:tsr}
\end{equation}
where $\mathbf{1}[\cdot]$ is the indicator function and equality is of the
entire sequence. A single substituted, inserted or dropped byte fails the whole
transmission. This is the quantity that matters to a link without
retransmission, and it is deliberately unforgiving: no partial credit is given
for a message that arrives nearly intact.

Character error rate measures how badly a failed transmission failed:
\begin{equation}
\mathrm{CER} = \frac{\sum_{i=1}^{N}
  d_{\mathrm{L}}\!\left(\hat{\mathbf{b}}^{(i)},\,\mathbf{b}^{(i)}\right)}
  {\sum_{i=1}^{N} L_i},
\label{eq:cer}
\end{equation}
where $d_{\mathrm{L}}$ is the Levenshtein distance. Three properties of
\eqref{eq:cer} are worth stating because they affect how the numbers should be
read. Levenshtein rather than Hamming distance is used because a decoder may
return a sequence of a different length from the one transmitted, so insertions
and deletions must be counted and not only substitutions. The ratio is a
micro-average, formed from the aggregate edit distance over the aggregate
transmitted length, so longer transmissions carry proportionally more weight
than shorter ones. And a decode that fails outright, returning no admissible
path, contributes $d_{\mathrm{L}} = L_i$ and is charged as total loss rather
than being excluded.

The term \emph{character} error rate is used in preference to \emph{byte} error
rate to avoid collision with the conventional meaning of BER as bit error rate;
since payloads are ASCII, one character is one byte throughout. The metric is
identical in the offline and live evaluations.

A reference level for achievable performance is also computed. A transmission
counts toward it if it was captured complete, or if it is missing exactly one
symbol and the capture carries the signature that symbol would leave: an
anomalous inter-pulse gap, or a merged pulse where a falling edge was missed. A
transmission missing exactly one symbol that left no such signature is
unrecoverable by any decoder operating on the detected pulse stream.

Two qualifications keep this from being a strict upper bound, and both make it
conservative. The signature test is applied only to the exactly-one-missing
class, so transmissions missing two or more symbols are counted as lost
outright even though some may carry enough timing evidence to be recovered;
on the held-out sets those number between 0 and 6 transmissions per rate.
Insertions are likewise not credited. The level is therefore best read as what
this decoder's mechanism can be expected to reach rather than as an
information-theoretic limit, and a decoder exceeding it would not be a
contradiction. It
is evaluated on the held-out transmissions only, the same population the
decoders are scored on, so that the three quantities plotted together in
Section~\ref{results} describe the same set of messages.

\subsection{Live Closed-Loop Operation}

In the live configuration the receiver streams pulses continuously, a decoder
process decodes each transmission as it completes and returns the decoded
message over the same serial link, and the receiver maintains a log of received
messages. Ground truth is polled from the beacon so that each decoded
transmission is scored against the message current at the time it was decoded.

Live runs were conducted at 28.6~bit/s and 114.3~bit/s, yielding 39 and 118
transmissions respectively. The decoder model was fitted from the sweep dataset
at the corresponding rate and was not refitted during the run; the live traffic
is therefore unseen data.

\subsection{Instrumentation Validation}

An artifact discovered during development is reported because it materially
affected an earlier measurement and illustrates a hazard of instrumenting an
embedded receiver.

In the original implementation the receiver wrote one line to its serial port
per detected pulse. The sampling loop halts for the duration of a write, and
the inter-symbol gap shrinks in proportion to the rate, so this dead time
consumes an increasing share of the interval available to detect the next
pulse. Taking a write to cost on the order of 1~ms, which is roughly four
iterations of the measured 262.4~$\mu$s detection loop but was not itself
measured, the dead time would occupy 10\% of the gap at 1$\times$ and 60\% at
6$\times$. Measured symbol loss under that implementation rose over the same
span from 0.21\% to 0.88\%, which is the signature of a fixed per-pulse dead
time rather than a degrading optical channel.

Buffering the pulse records and emitting them during the inter-message rest
reduced measured loss at every rate, though not uniformly: from 0.21\% to zero
at 1$\times$, 0.37\% to 0.32\% at 2$\times$, 0.59\% to 0.44\% at 3$\times$,
0.76\% to 0.37\% at 4$\times$, and 0.88\% to 0.52\% at 6$\times$. As a ratio
that is a factor of 1.15, 1.35, 2.04 and 1.71 at 2$\times$ through 6$\times$,
with loss eliminated outright at 1$\times$. The improvement is largest where
the dead time occupies the largest share of the inter-symbol gap, but it
reaches a full factor of two only at 4$\times$. All results in
Section~\ref{results} use the buffered implementation. The comparison is
reported because the earlier measurements were initially interpreted as channel
degradation, and the distinction was resolvable only by changing the
instrumentation and repeating the experiment.
